% Options for packages loaded elsewhere
\PassOptionsToPackage{unicode}{hyperref}
\PassOptionsToPackage{hyphens}{url}
\PassOptionsToPackage{dvipsnames,svgnames,x11names}{xcolor}
%
\documentclass[
]{article}

\usepackage{amsmath,amssymb}
\usepackage{iftex}
\ifPDFTeX
  \usepackage[T1]{fontenc}
  \usepackage[utf8]{inputenc}
  \usepackage{textcomp} % provide euro and other symbols
\else % if luatex or xetex
  \usepackage{unicode-math}
  \defaultfontfeatures{Scale=MatchLowercase}
  \defaultfontfeatures[\rmfamily]{Ligatures=TeX,Scale=1}
\fi
\usepackage{lmodern}
\ifPDFTeX\else  
    % xetex/luatex font selection
\fi
% Use upquote if available, for straight quotes in verbatim environments
\IfFileExists{upquote.sty}{\usepackage{upquote}}{}
\IfFileExists{microtype.sty}{% use microtype if available
  \usepackage[]{microtype}
  \UseMicrotypeSet[protrusion]{basicmath} % disable protrusion for tt fonts
}{}
\makeatletter
\@ifundefined{KOMAClassName}{% if non-KOMA class
  \IfFileExists{parskip.sty}{%
    \usepackage{parskip}
  }{% else
    \setlength{\parindent}{0pt}
    \setlength{\parskip}{6pt plus 2pt minus 1pt}}
}{% if KOMA class
  \KOMAoptions{parskip=half}}
\makeatother
\usepackage{xcolor}
\setlength{\emergencystretch}{3em} % prevent overfull lines
\setcounter{secnumdepth}{5}
% Make \paragraph and \subparagraph free-standing
\makeatletter
\ifx\paragraph\undefined\else
  \let\oldparagraph\paragraph
  \renewcommand{\paragraph}{
    \@ifstar
      \xxxParagraphStar
      \xxxParagraphNoStar
  }
  \newcommand{\xxxParagraphStar}[1]{\oldparagraph*{#1}\mbox{}}
  \newcommand{\xxxParagraphNoStar}[1]{\oldparagraph{#1}\mbox{}}
\fi
\ifx\subparagraph\undefined\else
  \let\oldsubparagraph\subparagraph
  \renewcommand{\subparagraph}{
    \@ifstar
      \xxxSubParagraphStar
      \xxxSubParagraphNoStar
  }
  \newcommand{\xxxSubParagraphStar}[1]{\oldsubparagraph*{#1}\mbox{}}
  \newcommand{\xxxSubParagraphNoStar}[1]{\oldsubparagraph{#1}\mbox{}}
\fi
\makeatother


\providecommand{\tightlist}{%
  \setlength{\itemsep}{0pt}\setlength{\parskip}{0pt}}\usepackage{longtable,booktabs,array}
\usepackage{calc} % for calculating minipage widths
% Correct order of tables after \paragraph or \subparagraph
\usepackage{etoolbox}
\makeatletter
\patchcmd\longtable{\par}{\if@noskipsec\mbox{}\fi\par}{}{}
\makeatother
% Allow footnotes in longtable head/foot
\IfFileExists{footnotehyper.sty}{\usepackage{footnotehyper}}{\usepackage{footnote}}
\makesavenoteenv{longtable}
\usepackage{graphicx}
\makeatletter
\def\maxwidth{\ifdim\Gin@nat@width>\linewidth\linewidth\else\Gin@nat@width\fi}
\def\maxheight{\ifdim\Gin@nat@height>\textheight\textheight\else\Gin@nat@height\fi}
\makeatother
% Scale images if necessary, so that they will not overflow the page
% margins by default, and it is still possible to overwrite the defaults
% using explicit options in \includegraphics[width, height, ...]{}
\setkeys{Gin}{width=\maxwidth,height=\maxheight,keepaspectratio}
% Set default figure placement to htbp
\makeatletter
\def\fps@figure{htbp}
\makeatother

\usepackage{booktabs}
\usepackage{longtable}
\usepackage{array}
\usepackage{multirow}
\usepackage{wrapfig}
\usepackage{float}
\usepackage{colortbl}
\usepackage{pdflscape}
\usepackage{tabu}
\usepackage{threeparttable}
\usepackage{threeparttablex}
\usepackage[normalem]{ulem}
\usepackage{makecell}
\usepackage{xcolor}
\usepackage{caption}
\usepackage{anyfontsize}
\usepackage{booktabs}
\usepackage{longtable}
\usepackage[a4paper, left=0.75in, right=1in, top=1in, bottom=1in]{geometry}
\makeatletter
\@ifpackageloaded{caption}{}{\usepackage{caption}}
\AtBeginDocument{%
\ifdefined\contentsname
  \renewcommand*\contentsname{Table of contents}
\else
  \newcommand\contentsname{Table of contents}
\fi
\ifdefined\listfigurename
  \renewcommand*\listfigurename{List of Figures}
\else
  \newcommand\listfigurename{List of Figures}
\fi
\ifdefined\listtablename
  \renewcommand*\listtablename{List of Tables}
\else
  \newcommand\listtablename{List of Tables}
\fi
\ifdefined\figurename
  \renewcommand*\figurename{Figure}
\else
  \newcommand\figurename{Figure}
\fi
\ifdefined\tablename
  \renewcommand*\tablename{Table}
\else
  \newcommand\tablename{Table}
\fi
}
\@ifpackageloaded{float}{}{\usepackage{float}}
\floatstyle{ruled}
\@ifundefined{c@chapter}{\newfloat{codelisting}{h}{lop}}{\newfloat{codelisting}{h}{lop}[chapter]}
\floatname{codelisting}{Listing}
\newcommand*\listoflistings{\listof{codelisting}{List of Listings}}
\makeatother
\makeatletter
\makeatother
\makeatletter
\@ifpackageloaded{caption}{}{\usepackage{caption}}
\@ifpackageloaded{subcaption}{}{\usepackage{subcaption}}
\makeatother

\ifLuaTeX
  \usepackage{selnolig}  % disable illegal ligatures
\fi
\usepackage{bookmark}

\IfFileExists{xurl.sty}{\usepackage{xurl}}{} % add URL line breaks if available
\urlstyle{same} % disable monospaced font for URLs
\hypersetup{
  pdftitle={Alpha Diversity},
  pdfauthor={Muluh},
  colorlinks=true,
  linkcolor={blue},
  filecolor={Maroon},
  citecolor={Blue},
  urlcolor={Blue},
  pdfcreator={LaTeX via pandoc}}


\title{Alpha Diversity}
\author{Muluh}
\date{}

\begin{document}
\maketitle


\section{Community diversity}\label{community-diversity}

\section{Group-wise comparisons}\label{group-wise-comparisons}

\begin{itemize}
\tightlist
\item
  \textbf{Index:} observed
\item
  \textbf{Group variable:} diet
\item
  \textbf{Multiple testing correction:} fdr
\item
  \textbf{Statistical test} Wilcoxon
\end{itemize}

\begin{center}
\includegraphics{alpha_observed_files/figure-pdf/visualize-plot-1.pdf}
\end{center}

\subsection{Table of the top findings}\label{table-of-the-top-findings}

\subsubsection{Across-group}\label{across-group}

\paragraph{Before treatment}\label{before-treatment}

\begin{itemize}
\tightlist
\item
  Multiple testing correction: \textbf{fdr}
\item
  Statistical test: \textbf{Wilcoxon}
\end{itemize}

\begin{longtable}[t]{lrrrrl}
\caption{Table of alpha diversity comparisons before Treatment}\\
\toprule
\textbf{group1} & \textbf{group2} & \textbf{mean1} & \textbf{mean2} & \textbf{log2FC} & \textbf{p.adj}\\
\midrule
\endfirsthead
\caption[]{Table of alpha diversity comparisons before Treatment \textit{(continued)}}\\
\toprule
\textbf{group1} & \textbf{group2} & \textbf{mean1} & \textbf{mean2} & \textbf{log2FC} & \textbf{p.adj}\\
\midrule
\endhead

\endfoot
\bottomrule
\endlastfoot
1 & 2 & 211.1 & 231.6 & 0.13 & 0.125\\*
\end{longtable}

\paragraph{After treatment}\label{after-treatment}

\begin{itemize}
\tightlist
\item
  Multiple testing correction: \textbf{fdr}
\item
  Statistical test: \textbf{Wilcoxon}
\end{itemize}

\begin{longtable}[t]{lrrrrl}
\caption{Table of alpha diversity comparisons after Treatment}\\
\toprule
\textbf{group1} & \textbf{group2} & \textbf{mean1} & \textbf{mean2} & \textbf{log2FC} & \textbf{p.adj}\\
\midrule
\endfirsthead
\caption[]{Table of alpha diversity comparisons after Treatment \textit{(continued)}}\\
\toprule
\textbf{group1} & \textbf{group2} & \textbf{mean1} & \textbf{mean2} & \textbf{log2FC} & \textbf{p.adj}\\
\midrule
\endhead

\endfoot
\bottomrule
\endlastfoot
1 & 2 & 207.1 & 233.3 & 0.17 & 0.053\\*
\end{longtable}

\begin{center}
\includegraphics{alpha_observed_files/figure-pdf/viz-div_change-1.pdf}
\end{center}

\subsubsection{Within-group}\label{within-group}

\paragraph{Paired treatment}\label{paired-treatment}

\begin{itemize}
\tightlist
\item
  Multiple testing correction: \textbf{paired} -none
\item
  Statistical test: \textbf{Wilcoxon}
\end{itemize}

\begin{longtable}[t]{lrrrr}
\caption{Table of alpha diversity paired between timepoint(1 and 2)}\\
\toprule
\textbf{diet} & \textbf{mean1} & \textbf{mean2} & \textbf{log2FC} & \textbf{p.value}\\
\midrule
\endfirsthead
\caption[]{Table of alpha diversity paired between timepoint(1 and 2) \textit{(continued)}}\\
\toprule
\textbf{diet} & \textbf{mean1} & \textbf{mean2} & \textbf{log2FC} & \textbf{p.value}\\
\midrule
\endhead

\endfoot
\bottomrule
\endlastfoot
1 & 211.1 & 207.1 & 0.04 & 0.128\\
2 & 231.6 & 233.3 & 0.13 & 0.390\\*
\end{longtable}

\begin{longtable}[t]{lrrrrl}
\caption{Table of interaction effects}\\
\toprule
\textbf{Term} & \textbf{Estimate} & \textbf{Std. Error} & \textbf{df} & \textbf{t value} & \textbf{Pr(>|t|)}\\
\midrule
\endfirsthead
\caption[]{Table of interaction effects \textit{(continued)}}\\
\toprule
\textbf{Term} & \textbf{Estimate} & \textbf{Std. Error} & \textbf{df} & \textbf{t value} & \textbf{Pr(>|t|)}\\
\midrule
\endhead

\endfoot
\bottomrule
\endlastfoot
(Intercept) & 7.673 & 0.085 & 88.101 & 90.739 & 0.000\\
diet2 & 0.115 & 0.119 & 88.097 & 0.971 & 0.334\\
time & -0.027 & 0.020 & 67.101 & -1.367 & 0.176\\
diet2:time & 0.037 & 0.028 & 67.101 & 1.302 & 0.197\\*
\end{longtable}

\section{Alternate analyses}\label{alternate-analyses}

\begin{itemize}
\item
  \textbf{Beta (Estimate)}: Shows the effect size and direction of a
  predictor on the outcome. Positive means an increase, negative a
  decrease.
\item
  \textbf{95\% CI}: Range where the true beta is likely to fall,
  reflecting precision.
\item
  \textbf{p-Value}: Tests if beta differs from zero; p\textless0.05p
  \textless{} 0.05p\textless0.05 suggests significance.
\end{itemize}

\begin{table}
\fontsize{12.0pt}{14.4pt}\selectfont
\begin{tabular*}{\linewidth}{@{\extracolsep{\fill}}lcccccc}
\toprule
 & \multicolumn{3}{c}{\textbf{before}} & \multicolumn{3}{c}{\textbf{after}} \\ 
\cmidrule(lr){2-4} \cmidrule(lr){5-7}
\textbf{Characteristic} & \textbf{Beta} & \textbf{95\% CI}\textsuperscript{\textit{1}} & \textbf{p-value} & \textbf{Beta} & \textbf{95\% CI}\textsuperscript{\textit{1}} & \textbf{p-value} \\ 
\midrule\addlinespace[2.5pt]
(Intercept) & 151 & 67, 236 & <0.001 & 144 & 59, 229 & 0.001 \\ 
meal\_group &  &  & 0.4 &  &  & 0.3 \\ 
    1 & — & — &  & — & — &  \\ 
    2 & 32 & -11, 75 & 0.14 & 28 & -16, 72 & 0.2 \\ 
    3 & 28 & -15, 71 & 0.2 & 30 & -14, 74 & 0.2 \\ 
    4 & 10 & -32, 53 & 0.6 & -3.9 & -47, 39 & 0.9 \\ 
age & 1.0 & -0.44, 2.5 & 0.2 & 1.2 & -0.25, 2.7 & 0.10 \\ 
\bottomrule
\end{tabular*}
\begin{minipage}{\linewidth}
\textsuperscript{\textit{1}}CI = Confidence Interval\\
\end{minipage}
\end{table}

\begin{table}
\fontsize{12.0pt}{14.4pt}\selectfont
\begin{tabular*}{\linewidth}{@{\extracolsep{\fill}}lccc}
\toprule
\textbf{Paired test between timepoints (1 and 2)} & \textbf{Beta} & \textbf{95\% CI}\textsuperscript{\textit{1}} & \textbf{p-value} \\ 
\midrule\addlinespace[2.5pt]
(Intercept) & 147 & 84, 210 & <0.001 \\ 
meal\_group &  &  & 0.6 \\ 
    1 & — & — &  \\ 
    2 & 28 & -16, 71 & 0.2 \\ 
    3 & 28 & -15, 71 & 0.2 \\ 
    4 & 2.1 & -28, 32 & 0.9 \\ 
age & 1.2 & 0.13, 2.2 & 0.025 \\ 
\bottomrule
\end{tabular*}
\begin{minipage}{\linewidth}
\textsuperscript{\textit{1}}CI = Confidence Interval\\
\end{minipage}
\end{table}

\section{Data Summary}\label{data-summary}

\begin{verbatim}
class: TreeSummarizedExperiment 
dim: 1700 140 
metadata(0):
assays(2): counts relabundance
rownames(1700): SGB4933 SGB4285 ... SGB15018 SGB4921
rowData names(9): kingdom phylum ... strain clade_name
colnames(140): AE1332 AK1371 ... TS2358 VK2336
colData names(15): sample id ... group observed_log2
reducedDimNames(0):
mainExpName: NULL
altExpNames(14): kingdom phylum ... KO metacyc
rowLinks: NULL
rowTree: NULL
colLinks: NULL
colTree: NULL
\end{verbatim}




\end{document}
